\documentclass[10pt,a4paper]{amsart}
\usepackage[margin=19mm]{geometry}
\usepackage{amsmath,amssymb,mathtools}
\usepackage[scr=rsfs,cal=euler]{mathalfa}
\usepackage{xfrac}
\usepackage[unicode,hidelinks,bookmarks=false]{hyperref}
\newcommand{\ie}{i.e.\ }
\newcommand{\cf}{cf.\ }
\newcommand{\resp}{respectively\ }
\newcommand{\Subsection}{\subsection}
\providecommand{\nobreakdash}{\nobreak\hbox{-}}
\setlength{\emergencystretch}{2em}
\multlinegap0pt
\usepackage{bm}
\usepackage{bbm}
\usepackage{dsfont}
\usepackage{xspace}
\usepackage{cancel}
\usepackage{mathtools}
\usepackage{amsthm}
\makeatletter
\@ifundefined{theorem}{\newtheorem{theorem}{Theorem}}{}
\@ifundefined{corollary}{\newtheorem{corollary}[theorem]{Corollary}}{}
\@ifundefined{proposition}{\newtheorem{proposition}[theorem]{Proposition}}{}
\makeatother
\newcommand{\HH}{\mathbb H}
\newcommand{\Tor}{\operatorname{Tor}}
\newcommand{\ZZ}{\mathbb Z}
\newcommand{\vol}{\operatorname{vol}}
\title[Cofinal towers with vanishing homology torsion]{Cofinal towers with vanishing homology torsion}
\author{Qilong Guo}
\date{Source: arXiv:2608.06601v1 (CC BY 4.0)}
\begin{document}
\maketitle

\noindent\emph{Source and licence.} Cofinal towers with vanishing homology torsion. Version arXiv:2608.06601v1 (CC BY 4.0), distributed under the Creative Commons Attribution 4.0 International licence, as recorded in the arXiv licence metadata for this exact version and shown on the abstract page https://arxiv.org/abs/2608.06601v1. Full rights notice: licenses/arxiv-2608.06601-CC-BY-4.0.txt.\par\smallskip

\noindent\textit{Editorial scope.} Counted unit: the Theorem whose source label is \texttt{thm:Girao} in the article (source0.tex lines 346--355, source label \texttt{thm:Girao}), delivered together with the article's own complete original proof of it (source0.tex lines 359--397, one \texttt{proof} environment of 39 lines). No mathematics is written, repaired or re-derived: statement and proof are the article's own text line for line. The proof is detached from the statement in the source: the article states the result at lines 346--355 and prints its proof at lines 359--397, and the proof environment's own header names the statement, so the association is the article's own. Required context retained uncounted: thm:main (lines 86--102); eq:H1-torsion-free (lines 257--259); prop:compatibility (lines 303--312); cor:double-cover (lines 324--329). Every definition, display and label of those lines is retained unchanged. Omitted from this package and not redistributed: 1--85, 103--256, 260--302, 313--323, 330--345, 356--358, 398--462. Nothing inside a retained excerpt is omitted or altered: every retained body is delivered exactly as the article's lines are written, so no omission receipt of any kind accompanies this package. Citations: the retained excerpts carry no citation command at all, so no bibliography entry of the article is transcribed and no reference is redistributed. Reference adaptation: none was needed and none was made; every reference inside the delivered unit resolves inside it, so no unresolved reference or citation is produced. Printed slips kept literal: no printed word, symbol, hypothesis or number of the counted statement or of its proof was altered. Layout and class: the article's own class is \texttt{amsart} with packages including fontenc, lmodern, microtype, amsmath, amssymb, amsthm, mathtools, geometry, hyperref; the wrapper is the lane's pinned \texttt{amsart} editorial wrapper at 10pt on A4, which restates the theorem-like environments the retained text needs and adds the article's own macro definitions that the retained text needs (\texttt{\textbackslash HH}, \texttt{\textbackslash Tor}, \texttt{\textbackslash ZZ}, \texttt{\textbackslash vol}), with extra wrapper packages (bm, bbm, dsfont, xspace, cancel, mathtools). The delivered numbering is the unit's own. Assets: no retained span contains a figure environment, a graphics-inclusion command, a PDF-page inclusion, a table or a TikZ picture; no image file is copied into this package, the package declares no asset dependency and no image file is redistributed with it. Licence: version arXiv:2608.06601v1 of this article is distributed under the Creative Commons Attribution 4.0 International licence, as recorded in the arXiv licence metadata for this exact version and shown on the abstract page https://arxiv.org/abs/2608.06601v1. A full rights notice is delivered at licenses/arxiv-2608.06601-CC-BY-4.0.txt. Characters: every retained excerpt is pure ASCII, so the delivered engine, XeLaTeX with its default Latin Modern text font, reports no missing character.
\begin{theorem}\label{thm:main}
Let $P_0\subset\HH^3$ be an ideal right-angled polyhedron, and let $M_0=N(P_0)$ be its checkerboard manifold.  There is a cofinal tower
\begin{equation}\label{eq:main-tower}
  M_0 \longleftarrow M_1 \longleftarrow M_2 \longleftarrow \cdots
\end{equation}
such that every bonding map $M_{n+1}\to M_n$ is a regular double cover and every $M_n$ is a hyperbolic link complement in $S^3$.  Moreover,
\[
  \deg(M_n\to M_0)=2^n,
  \qquad
  \Tor H_1(M_n;\ZZ)=0
\]
for every $n$.  In particular,
\[
  \frac{\log |\Tor H_1(M_n;\ZZ)|}{\vol(M_n)}=0
\]
at every level.
\end{theorem}

\begin{equation}\label{eq:H1-torsion-free}
  \Tor H_1(N(P);\ZZ)=0.
\end{equation}

\begin{proposition}\label{prop:compatibility}
Under the inclusion $W(P')<W(P)$,
\begin{equation}\label{eq:compatibility}
  \epsilon_{P'}=\left.\epsilon_P\right|_{W(P')}.
\end{equation}
Consequently,
\begin{equation}\label{eq:kernel-pullback}
  \Gamma(P')=W(P')\cap\Gamma(P).
\end{equation}
\end{proposition}

\begin{corollary}\label{cor:double-cover}
Reflection doubling induces a regular double cover
\begin{equation}\label{eq:double-cover}
  N(P')\longrightarrow N(P).
\end{equation}
\end{corollary}

\begin{theorem}[Gir\~ao]\label{thm:Girao}
For every ideal right-angled polyhedron $P_0$, there is a sequence
\begin{equation}\label{eq:polyhedron-sequence}
  P_0,P_1,P_2,\ldots
\end{equation}
in which $P_{n+1}$ is a reflection double of $P_n$ and
\begin{equation}\label{eq:W-cofinal}
  \bigcap_{n\geq0}W(P_n)=\{1\}.
\end{equation}
\end{theorem}

\begin{proof}[Proof of Theorem~\ref{thm:main}]
Choose $P_n$ as in Theorem~\ref{thm:Girao}, and set
\[
  W_n=W(P_n),
  \qquad
  \Gamma_n=\Gamma(P_n),
  \qquad
  M_n=N(P_n).
\]
Proposition~\ref{prop:compatibility} and Corollary~\ref{cor:double-cover} give
\[
  \Gamma_{n+1}=W_{n+1}\cap\Gamma_n,
  \qquad
  [\Gamma_n:\Gamma_{n+1}]=2.
\]
Thus $M_{n+1}\to M_n$ is a regular double cover and
\begin{equation}\label{eq:degree}
  \deg(M_n\to M_0)=2^n.
\end{equation}

Since $\Gamma_n\subset W_n$,
\[
  \bigcap_{n\geq0}\Gamma_n
  \subset
  \bigcap_{n\geq0}W_n
  =\{1\},
\]
so the tower is cofinal.

By~\eqref{eq:H1-torsion-free},
\[
  \Tor H_1(M_n;\ZZ)=0
\]
for every $n$, and therefore
\[
  \frac{\log |\Tor H_1(M_n;\ZZ)|}{\vol(M_n)}=0.
\]

\end{proof}

\end{document}
